\documentclass[12pt]{article}

\input{~/School/Notes/header.tex}

\title{Домашнее задание по алгебре}
\author{Денисов Илья}
\date{}

\begin{document}
	\maketitle

	\section*{Задача 1}

	\noindent
	\textbf{Условие.} Многочлен $P(x) = f(x^2) + x \cdot g(x^2)$ делится на $x^2 - 2$. Докажите, что $f$ имеет вещественный корень.

	\begin{proof}
		Так как $P$ делится на $x^2 - 2$, найдётся многочлен $Q$ такой, что
		\[
			P(x) = (x^2 - 2) \cdot Q(x)
		\]
		Числа $\sqrt{2}$ и $-\sqrt{2}$ обращают $x^2 - 2$ в ноль, поэтому
		\[
			P(\sqrt{2}) = 0, \qquad P(-\sqrt{2}) = 0
		\]
		Выпишем эти два равенства явно. Заметим, что $(\sqrt{2})^2 = (-\sqrt{2})^2 = 2$. Поэтому в обоих случаях внутри $f$ и $g$ стоит число $2$:
		\begin{align*}
			P(\sqrt{2}) &= f\bigl((\sqrt{2})^2\bigr) + \sqrt{2} \cdot g\bigl((\sqrt{2})^2\bigr) = f(2) + \sqrt{2}\, g(2) = 0 \\
			P(-\sqrt{2}) &= f\bigl((-\sqrt{2})^2\bigr) - \sqrt{2} \cdot g\bigl((-\sqrt{2})^2\bigr) = f(2) - \sqrt{2}\, g(2) = 0
		\end{align*}
		Сложим два равенства. Слагаемые $\pm\sqrt{2}\, g(2)$ взаимно уничтожатся:
		\[
			2 f(2) = 0 \quad \Longrightarrow \quad f(2) = 0
		\]
		Значит, вещественное число $2$ является корнем многочлена $f$.
	\end{proof}

	\section*{Задача 2}

	\noindent
	\textbf{Условие.} а) Найдите многочлен $g$ от четырёх переменных $x_1, x_2, x_3, x_4$, из которого перестановками переменных получаются ровно три многочлена $g_1 = g$, $g_2$, $g_3$. б) Выразите сумму, сумму попарных произведений и произведение $g_1, g_2, g_3$ через сумму, сумму попарных произведений и т.\,д. самих переменных $x_1, x_2, x_3, x_4$.

	\subsection*{а)}

	\noindent
	\textbf{Ответ.} Подходит многочлен
	\[
		g = x_1x_2 + x_3x_4
	\]
	Три получающихся многочлена:
	\[
		g_1 = x_1x_2 + x_3x_4, \qquad g_2 = x_1x_3 + x_2x_4, \qquad g_3 = x_1x_4 + x_2x_3
	\]

	\begin{proof}[Обоснование]
		Любая перестановка переменных превращает $g$ в многочлен вида $x_ax_b + x_cx_d$, где $\{a, b, c, d\} = \{1, 2, 3, 4\}$. Такой многочлен зависит только от того, как множество индексов $\{1, 2, 3, 4\}$ разбито на две неупорядоченные пары $\{a, b\}$ и $\{c, d\}$. Действительно, порядок внутри пары не важен, потому что $x_ax_b = x_bx_a$. Порядок самих пар тоже не важен, потому что слагаемые можно поменять местами.

		Разбиений четырёх элементов на две пары ровно три:
		\[
			\{1, 2\} \,|\, \{3, 4\}, \qquad \{1, 3\} \,|\, \{2, 4\}, \qquad \{1, 4\} \,|\, \{2, 3\}
		\]
		Им отвечают $g_1$, $g_2$, $g_3$. Эти многочлены попарно различны, потому что, например, одночлен $x_1x_2$ входит только в $g_1$. Значит, получается ровно три многочлена.
	\end{proof}

	\subsection*{б)}

	Обозначим элементарные симметрические многочлены от $x_1, x_2, x_3, x_4$:
	\begin{align*}
		\alpha_1 &= x_1 + x_2 + x_3 + x_4 \\
		\alpha_2 &= x_1x_2 + x_1x_3 + x_1x_4 + x_2x_3 + x_2x_4 + x_3x_4 \\
		\alpha_3 &= x_1x_2x_3 + x_1x_2x_4 + x_1x_3x_4 + x_2x_3x_4 \\
		\alpha_4 &= x_1x_2x_3x_4
	\end{align*}

	\medskip
	\noindent
	\textbf{1) Сумма.} В сумме $g_1 + g_2 + g_3$ каждое из шести попарных произведений $x_ix_j$ встречается ровно один раз, поэтому
	\[
		\boxed{g_1 + g_2 + g_3 = \alpha_2}
	\]

	\medskip
	\noindent
	\textbf{2) Сумма попарных произведений.} Перемножим многочлены попарно:
	\begin{align*}
		g_1g_2 &= (x_1x_2 + x_3x_4)(x_1x_3 + x_2x_4) = x_1^2x_2x_3 + x_1x_2^2x_4 + x_1x_3^2x_4 + x_2x_3x_4^2 \\
		g_1g_3 &= (x_1x_2 + x_3x_4)(x_1x_4 + x_2x_3) = x_1^2x_2x_4 + x_1x_2^2x_3 + x_1x_3x_4^2 + x_2x_3^2x_4 \\
		g_2g_3 &= (x_1x_3 + x_2x_4)(x_1x_4 + x_2x_3) = x_1^2x_3x_4 + x_1x_2x_3^2 + x_1x_2x_4^2 + x_2^2x_3x_4
	\end{align*}
	Сложим и сгруппируем слагаемые по тому, какая переменная стоит в квадрате:
	\begin{align*}
		g_1g_2 + g_1g_3 + g_2g_3 = {}& x_1^2(x_2x_3 + x_2x_4 + x_3x_4) + x_2^2(x_1x_3 + x_1x_4 + x_3x_4) \\
		&+ x_3^2(x_1x_2 + x_1x_4 + x_2x_4) + x_4^2(x_1x_2 + x_1x_3 + x_2x_3)
	\end{align*}
	Получилось $12$ слагаемых вида $x_i^2x_jx_k$: для каждого $i$ берутся все три пары $\{j, k\}$ из оставшихся индексов.

	Теперь посчитаем $\alpha_1\alpha_3$. Умножим $x_i$ на каждое из четырёх слагаемых $\alpha_3$. Если $x_i$ входит в тройку, получится $x_i^2x_jx_k$. Если не входит, получится $x_1x_2x_3x_4 = \alpha_4$. Для каждой из четырёх переменных есть ровно одна тройка, в которую она не входит, поэтому слагаемое $\alpha_4$ встретится $4$ раза. Итак,
	\[
		\alpha_1\alpha_3 = \underbrace{\sum_{i} x_i^2 \sum_{\substack{j < k \\ j, k \neq i}} x_jx_k}_{\text{те же 12 слагаемых}} + 4\alpha_4
	\]
	Сравнивая с предыдущим вычислением, получаем
	\[
		\boxed{g_1g_2 + g_1g_3 + g_2g_3 = \alpha_1\alpha_3 - 4\alpha_4}
	\]

	\medskip
	\noindent
	\textbf{3) Произведение.} Умножим найденное выше $g_1g_2$ на $g_3 = x_1x_4 + x_2x_3$:
	\begin{align*}
		g_1g_2g_3 = {}& \bigl(x_1^2x_2x_3 + x_1x_2^2x_4 + x_1x_3^2x_4 + x_2x_3x_4^2\bigr)(x_1x_4 + x_2x_3) \\
		= {}& \underbrace{x_1^3x_2x_3x_4 + x_1x_2^3x_3x_4 + x_1x_2x_3^3x_4 + x_1x_2x_3x_4^3}_{= \alpha_4 (x_1^2 + x_2^2 + x_3^2 + x_4^2)} \\
		&+ \underbrace{x_1^2x_2^2x_3^2 + x_1^2x_2^2x_4^2 + x_1^2x_3^2x_4^2 + x_2^2x_3^2x_4^2}_{\text{квадраты четырёх слагаемых } \alpha_3}
	\end{align*}
	Разберём каждую часть.

	\begin{itemize}
		\item Сумма квадратов: $x_1^2 + x_2^2 + x_3^2 + x_4^2 = \alpha_1^2 - 2\alpha_2$, потому что в $\alpha_1^2$ квадраты входят по одному разу, а каждое произведение $x_ix_j$ ($i \neq j$) входит дважды.
		\item Сумма квадратов слагаемых $\alpha_3$. Так как $\alpha_3^2$ равно сумме квадратов слагаемых плюс удвоенная сумма всех попарных произведений различных слагаемых, то
		\[
			x_1^2x_2^2x_3^2 + \dots + x_2^2x_3^2x_4^2 = \alpha_3^2 - 2S
		\]
		где $S$~--- сумма произведений двух различных троек. Две различные тройки имеют вид $\{i, j, k\}$ и $\{i, j, l\}$, и их произведение равно $x_ix_j \cdot x_ix_jx_kx_l = x_ix_j\alpha_4$. Каждая пара $\{i, j\}$ встречается ровно один раз, поэтому $S = \alpha_4(x_1x_2 + x_1x_3 + \dots + x_3x_4) = \alpha_4\alpha_2$.
	\end{itemize}
	Собираем всё вместе:
	\[
		g_1g_2g_3 = \alpha_4(\alpha_1^2 - 2\alpha_2) + \alpha_3^2 - 2\alpha_2\alpha_4 = \alpha_1^2\alpha_4 + \alpha_3^2 - 4\alpha_2\alpha_4
	\]
	\[
		\boxed{g_1g_2g_3 = \alpha_1^2\alpha_4 + \alpha_3^2 - 4\alpha_2\alpha_4}
	\]

	%==================================================
	\section*{Задача 3}

	\noindent
	\textbf{Условие.} а) Укажите многочлен $6$-й степени с целыми коэффициентами, корнем которого является корень $7$-й степени из $1$ (но не единица). б) Укажите многочлен $3$-й степени с целыми коэффициентами, корнем которого является $\cos\dfrac{2\pi}{7}$. в) Вычислите $\dfrac{1}{\cos\frac{2\pi}{7}} + \dfrac{1}{\cos\frac{4\pi}{7}} + \dfrac{1}{\cos\frac{6\pi}{7}}$.
	
	\subsection*{а)}

	Пусть $\alpha = \cos\dfrac{2\pi}{7} + i\sin\dfrac{2\pi}{7}$. Тогда $\alpha^7 = 1$ и $\alpha \neq 1$. Разложим многочлен $x^7 - 1$ на множители:
	\[
		x^7 - 1 = (x - 1)(x^6 + x^5 + x^4 + x^3 + x^2 + x + 1)
	\]
	Подставим $x = \alpha$: $0 = \alpha^7 - 1 = (\alpha - 1)\,\Phi(\alpha)$, где $\Phi(x) = x^6 + x^5 + x^4 + x^3 + x^2 + x + 1$. Так как $\alpha - 1 \neq 0$, то $\Phi(\alpha) = 0$.

	\medskip
	\noindent
	\textbf{Ответ.} $\Phi(x) = x^6 + x^5 + x^4 + x^3 + x^2 + x + 1$.

	\subsection*{б)}

	Пусть $\alpha = e^{2\pi i/7}$ и $t = \alpha + \alpha^{-1}$. Так как $|\alpha| = 1$, то $\alpha^{-1} = \overline{\alpha}$, и
	\[
		t = \alpha + \overline{\alpha} = 2\operatorname{Re}\alpha = 2\cos\frac{2\pi}{7}
	\]
	Разделим равенство $\Phi(\alpha) = 0$ на $\alpha^3$ (это можно сделать, так как $\alpha \neq 0$), а затем сгруппируем слагаемые парами:
	\[
		\alpha^3 + \alpha^2 + \alpha + 1 + \alpha^{-1} + \alpha^{-2} + \alpha^{-3} = 0
	\]
	\[
		(\alpha^3 + \alpha^{-3}) + (\alpha^2 + \alpha^{-2}) + (\alpha + \alpha^{-1}) + 1 = 0
	\]
	Выразим каждую скобку через $t$:
	\begin{align*}
		t^2 &= \alpha^2 + 2 + \alpha^{-2} & &\Longrightarrow & \alpha^2 + \alpha^{-2} &= t^2 - 2 \\
		t^3 &= \alpha^3 + 3\alpha + 3\alpha^{-1} + \alpha^{-3} & &\Longrightarrow & \alpha^3 + \alpha^{-3} &= t^3 - 3t
	\end{align*}
	Подставляем:
	\[
		(t^3 - 3t) + (t^2 - 2) + t + 1 = 0 \quad \Longleftrightarrow \quad t^3 + t^2 - 2t - 1 = 0
	\]
	Теперь вернёмся к косинусу: $t = 2c$, где $c = \cos\dfrac{2\pi}{7}$. Тогда $8c^3 + 4c^2 - 4c - 1 = 0$.

	\medskip
	\noindent
	\textbf{Ответ.} $8x^3 + 4x^2 - 4x - 1$.

	\subsection*{в)}

	Обозначим $c_k = \cos\dfrac{2\pi k}{7}$ для $k = 1, 2, 3$. Тогда нужная сумма равна $\dfrac{1}{c_1} + \dfrac{1}{c_2} + \dfrac{1}{c_3}$.

	\medskip
	\noindent
	\textbf{Шаг 1: все три $c_k$~--- корни одного многочлена.} В пункте б) мы использовали только то, что $\Phi(\alpha) = 0$. Число $\alpha^k$ ($k = 1, 2, 3$) тоже является корнем $7$-й степени из $1$, отличным от $1$, поэтому $\Phi(\alpha^k) = 0$. Повторяя вычисление с $\alpha^k$ вместо $\alpha$, получаем, что $c_k = \cos\dfrac{2\pi k}{7}$ является корнем многочлена $8x^3 + 4x^2 - 4x - 1$.

	Числа $c_1, c_2, c_3$ попарно различны. Углы $\dfrac{2\pi}{7} < \dfrac{4\pi}{7} < \dfrac{6\pi}{7}$ лежат в промежутке $(0, \pi)$, где косинус строго убывает. У кубического многочлена не больше трёх корней, значит, $c_1, c_2, c_3$~--- это \emph{все} его корни.

	\medskip
	\noindent
	\textbf{Шаг 2: теорема Виета.} Для многочлена $ax^3 + bx^2 + cx + d$ с корнями $r_1, r_2, r_3$ верно
	\[
		r_1 + r_2 + r_3 = -\frac{b}{a}, \qquad r_1r_2 + r_1r_3 + r_2r_3 = \frac{c}{a}, \qquad r_1r_2r_3 = -\frac{d}{a}
	\]
	Для $8x^3 + 4x^2 - 4x - 1$ имеем $a = 8$, $b = 4$, $c = -4$, $d = -1$, поэтому
	\[
		c_1 + c_2 + c_3 = -\frac{1}{2}, \qquad c_1c_2 + c_1c_3 + c_2c_3 = -\frac{1}{2}, \qquad c_1c_2c_3 = \frac{1}{8}
	\]
	В частности, произведение не равно нулю, поэтому все $c_k \neq 0$ и дроби $\dfrac{1}{c_k}$ определены.

	\medskip
	\noindent
	\textbf{Шаг 3: приводим к общему знаменателю.}
	\[
		\frac{1}{c_1} + \frac{1}{c_2} + \frac{1}{c_3} = \frac{c_2c_3 + c_1c_3 + c_1c_2}{c_1c_2c_3} = \frac{-\frac{1}{2}}{\frac{1}{8}} = -4
	\]

	\medskip
	\noindent
	\textbf{Ответ.} $-4$.

	%==================================================
	\section*{Задача 4}

	\noindent
	\textbf{Условие.} У неприводимого над целыми числами многочлена со старшим коэффициентом $1$ есть корень, по модулю равный $1$. Докажите, что многочлен возвратный (то есть переходит сам в себя при замене $x$ на $1/x$ и домножении на степень $x$) и что его степень чётна.

	%==================================================
	\section*{Задача 5}

	\noindent
	\textbf{Условие.} Избавьтесь от иррациональности в знаменателе в выражении $\dfrac{1}{1 + \sqrt[3]{2} + 2\sqrt[3]{4}}$.

	\medskip
	\noindent
	\textbf{Идея.} Обозначим $a = \sqrt[3]{2}$. Тогда $a^3 = 2$ и $\sqrt[3]{4} = a^2$, а знаменатель равен $D = 1 + a + 2a^2$. Будем искать число $N = p + qa + ra^2$ с рациональными $p, q, r$ такое, что $D \cdot N = 1$. Тогда $\dfrac{1}{D} = N$.

	\medskip
	
	Перемножаем. Раскроем скобки и воспользуемся тем, что $a^3 = 2$ и $a^4 = a \cdot a^3 = 2a$:
	\begin{align*}
		D \cdot N &= (1 + a + 2a^2)(p + qa + ra^2) \\
		&= p + qa + ra^2 + pa + qa^2 + ra^3 + 2pa^2 + 2qa^3 + 2ra^4 \\
		&= p + qa + ra^2 + pa + qa^2 + 2r + 2pa^2 + 4q + 4ra
	\end{align*}
	Соберём слагаемые при $1$, $a$ и $a^2$:
	\[
		D \cdot N = (p + 4q + 2r) + (p + q + 4r)\, a + (2p + q + r)\, a^2
	\]

	Мы хотим, чтобы $D \cdot N = 1 + 0 \cdot a + 0 \cdot a^2$.
	\[
		\begin{cases}
			p + 4q + 2r = 1 \\
			p + q + 4r = 0 \\
			2p + q + r = 0
		\end{cases}
	\]

	Решаем систему. Вычтем второе уравнение из третьего:
	\[
		(2p + q + r) - (p + q + 4r) = p - 3r = 0 \quad \Longrightarrow \quad p = 3r
	\]
	Подставим $p = 3r$ в третье уравнение: $6r + q + r = 0$, откуда $q = -7r$. Подставим оба выражения в первое уравнение:
	\[
		3r - 28r + 2r = -23r = 1 \quad \Longrightarrow \quad r = -\frac{1}{23}, \quad p = -\frac{3}{23}, \quad q = \frac{7}{23}
	\]

	Значит, $D \cdot \dfrac{-3 + 7a - a^2}{23} = 1$.

	\medskip
	\noindent
	\textbf{Ответ.}
	\[
		\boxed{\frac{1}{1 + \sqrt[3]{2} + 2\sqrt[3]{4}} = \frac{7\sqrt[3]{2} - \sqrt[3]{4} - 3}{23}}
	\]
\end{document}
